\subsection{\texorpdfstring{映射 \(\Theta_{a}\) 的构造}{映射 \textbackslash Theta\_\{a\} 的构造}}

设 \(\boldsymbol{a}=(a_{1},\ldots,a_{n})\in\mathrm{A}^{n}\)，并设 U 是 \(\mathrm{Sp}^{n}(\boldsymbol{a})\) 的一个开邻域。设 h 是一个无穷次可微函数，在 \(\mathrm{Sp}^{n}(\boldsymbol{a})\) 的一个邻域中等于 1，且 h 的支集是紧的并包含于 U（I, p.~52, 引理 1）。依 I, p.~54, 引理 3，存在无穷次可微函数 \((u_{1},\ldots,u_{n})\)（\(C^{n}\) 上的、取值于 A 中的），使得序列 \((h,u_{1},\ldots,u_{n})\) 适配于 a。设 \(\omega\) 是相伴的微分形式；它具有包含于 U 中的紧支集（I, p.~54, 引理 4）。存在一个无穷次可微函数 \(\psi\)（在 U 中具有紧支集且取值于 A 中），使得

\[
\omega = \psi d x _ {1} \wedge d y _ {1} \wedge \dots \wedge d x _ {n} \wedge d y _ {n}
\]

其中 \(x_{j} + iy_{j}\) 是 \(C^{n}\)（被等同于 \(R^{2n}\)）上的坐标函数。设 \(\mu\) 是 \(R^{2n}\) 上的 Lebesgue 测度。

设 \(f \in \mathcal{C}(\mathrm{U};\mathrm{A})\)。微分形式 \(f\omega|\mathrm{U}\)（\(\mathrm{U}\) 上的）是连续的且具有紧支集。与这个微分形式相伴的向量测度（VAR, R2, 10.4.3 与 10.4.4）是向量测度 \(f\psi \cdot \mu\)；它是一个以 \(\mu\) 为底的测度（INT, VI, §2, no. 4, def. 4），它的支集是紧的且包含于 \(\omega\) 的支集。这个测度相对于 A 的范数是有界的（INT, VI, §2, no. 4, prop. 8）；我们用 \(\| f\omega\|\) 表示与之相伴的 \(\mathbf{C}^n\) 上的正测度（INT, VI, §2, no. 3, def. 3）。

依 INT, VI, §2, no. 4, prop. 8, b)，我们有

\[
\| f \omega \| = \| f \psi \cdot \mu \| = \| f \psi \| \cdot \mu = \| f \| \| \omega \|.
\]

特别地，微分形式 \(f\omega\) 在 U 上的积分满足

\[
\left\| \int_ {\mathrm{U}} f \omega \right\| \leqslant \int_ {\mathrm{U}} \| f \| \| \omega \|
\]

（INT, VI, §2, no. 3, prop. 5）。

引理 7. --- 对每个函数 \(f \in \mathcal{O}(\mathrm{U};\mathrm{A})\)，积分 \(\int_{\mathrm{U}} f\omega\) 是 A 的一个元素，它只依赖于 \(\boldsymbol{a}\) 以及 \(f\) 在 \(\mathrm{Sp}^{n}(\boldsymbol{a})\) 的一个邻域中的芽。它满足不等式

\[
\left\| \int_ {\mathrm{U}} f \omega \right\| \leqslant \left(\int_ {\mathrm{U}} \| \omega \|\right) \sup _ {z \in \operatorname{Supp} (h)} \| f (z) \|,\tag{4}
\]

并且

\[
\int_ {\mathrm{U}} (a f) \omega = a \int_ {\mathrm{U}} f \omega
\]

对一切 \(a \in \mathbf{A}\)。

我们在上面已经看到，积分对 \(f \in \mathcal{C}(\mathrm{U};\mathrm{A})\) 有定义，并且满足

\[
\left\| \int_ {\mathrm{U}} f \omega \right\| \leqslant \int_ {\mathrm{U}} \| f \| \| \omega \| \leqslant \left(\int_ {\mathrm{U}} \| \omega \|\right) \sup _ {z \in \operatorname{Supp} (h)} \| f (z) \|.
\]

此外，对每个 \(a \in \mathrm{A}\) 与每个 \(f \in \mathcal{C}(\mathrm{U};\mathrm{A})\)，我们有

\[
\int_ {\mathrm{U}} (a f) \omega = a \int_ {\mathrm{U}} f \omega
\]

（INT, VI, §2, no. 2, prop. 2，应用于乘以 a）。

设 \((h', u_{1}', \ldots, u_{n}')\) 是一个适配于 a 的序列，满足 \(\operatorname{Supp}(h') \subset \operatorname{U}\)。设 \(\omega'\) 是相伴的微分形式。依 I, p.~57, 引理 6，存在一个 n-1 次微分形式 \(\psi\)（\(C^{n}\) 上的）（以 A 为系数且支包含于 \(\operatorname{Supp}(h) \cup \operatorname{Supp}(h') \subset \operatorname{U}\) 中），使得

\[
\omega - \omega^ {\prime} = d (\psi \wedge d z _ {1} \wedge \dots \wedge d z _ {n}).
\]

设 \(f \in \mathcal{O}(\mathrm{U};\mathrm{A})\)。由于映射 \(f\) 是全纯的，我们有

\[
d f = \sum_ {i = 1} ^ {n} \frac {\partial f}{\partial z _ {i}} d z _ {i},
\]

（VAR, R2, p.~24, 8.8.9），因而

\[
f \left(\omega - \omega^ {\prime}\right) = f d \left(\psi \wedge d z _ {1} \wedge \dots \wedge d z _ {n}\right) = d \left(f \psi \wedge d z _ {1} \wedge \dots \wedge d z _ {n}\right).
\]

依 Stokes 公式（VAR, R2, p.~48, 11.2.3），我们于是有

\[
\int_ {\mathrm{U}} f (\omega - \omega^ {\prime}) = 0.
\]

于是元素 \(\int_{\mathrm{U}} f\omega\) 不依赖于序列 \((h, u_1, \ldots, u_n)\) 的选取。为完成证明，我们来证 \(\int_{\mathrm{U}} f\omega\) 只依赖于 \(f\) 在 \(\mathrm{Sp}^n(\boldsymbol{a})\) 的一个邻域中的芽。设 U 与 \(\mathrm{U}'\) 是 \(\mathrm{Sp}^n(\boldsymbol{a})\) 的两个开邻域。设 \(f \in \mathcal{O}(\mathrm{U}; \mathrm{A})\) 与 \(f' \in \mathcal{O}(\mathrm{U}', \mathrm{A})\)，使得 \(f\) 与 \(f'\) 在一个开邻域 \(\mathrm{U}''\)（\(\mathrm{Sp}^n(\boldsymbol{a})\) 的）上相等。存在一个映射 \(h\)（从 \(\mathbf{C}^n\) 到 \(\mathbf{C}\) 中的、无穷次可微的），在 \(\mathrm{Sp}^n(\boldsymbol{a})\) 的一个邻域中等于 1，且具有包含于 \(\mathrm{U}''\) 中的紧支集（I, p.~52, 引理 1），并且存在 \((u_1, \ldots, u_n)\)，使得序列 \((h, u_1, \ldots, u_n)\) 适配于 \(\boldsymbol{a}\)（I, p.~54, 引理 3）。设 \(\omega\) 是相伴的微分形式。由于 \(\mathrm{Supp}(\omega) \subset \mathrm{Supp}(h) \subset \mathrm{U}''\)（I, p.~54, 引理 4, b)），我们有

\[
\int_ {\mathrm{U}} f \omega = \int_ {\mathrm {U^ {\prime\prime}}} f \omega = \int_ {\mathrm {U^ {\prime\prime}}} f ^ {\prime} \omega = \int_ {\mathrm {U^ {\prime}}} f ^ {\prime} \omega ,
\]

这就完成了证明。

这个引理表明，存在唯一的 A-线性映射 \(\Theta_{\boldsymbol{a}}\)（从 \(\mathcal{O}(\mathrm{Sp}^{n}(\boldsymbol{a});\mathrm{A})\) 到 A 中的），使得

\[
\Theta_ {\boldsymbol {a}} (f) = \frac {n !}{(2 i \pi) ^ {n}} \int_ {\mathrm{U}} \widetilde {f} \omega\tag{5}
\]

它对每个开集 U 与每个代表元 \(\tilde{f} \in \mathcal{O}(\mathrm{U};\mathrm{A})\)（某个芽 \(f \in \mathcal{O}(\mathrm{Sp}^n (\boldsymbol {a});\mathrm{A})\) 的）成立。线性映射 \(\Theta_{\boldsymbol{a}}\) 依不等式 (4) 以及 EVT, II, p.~29, prop. 5 是连续的。
% semantic-fragments: legacy-input-seam
\relax{}
